\chapter{Angular Momentum, Spin, and Statistics}

\section{From Ordinary Rotation to Intrinsic Spin}

Angular momentum is the formal subject, but spin is the real target. Spin is one of the simplest and most unintuitive properties of elementary particles, and its structure follows from the mathematics of angular momentum.

Angular momentum is a vector. Its direction lies along the axis of rotation. There are two possible directions along that axis, so a convention is needed to choose between them: curl the fingers of the right hand in the sense of rotation, and the thumb points in the direction of the angular-momentum vector. This is a definition, not a theorem.

The magnitude depends on how much mass is present, how that mass is distributed, the size of the system, and its angular velocity. The mass distribution and size are collected into the moment of inertia. A large object and a small object can therefore carry the same angular momentum while rotating at very different rates and with very different energies.

An ordinary composite object can carry two conceptually different kinds of angular momentum. A cup moving in a circle around an external axis has angular momentum because its center of mass is moving around that axis, even if the cup is not turning about any axis through itself. This is orbital angular momentum. Spin means rotation about an internal axis.

The clean definition uses the center-of-mass frame. Go to the frame in which the center of mass is at rest, so the total ordinary momentum vanishes. Any angular momentum that remains in that frame is spin angular momentum. The object need not be motionless internally; only its center of mass is at rest.

For a basketball this distinction causes no difficulty. A basketball at rest and the same basketball set rotating are plainly recognized as the same object in two different states. Quantum mechanics complicates this picture because angular momentum is discrete. There is no continuous sequence of allowed angular momenta joining a state of zero angular momentum to the first rotating state. One can therefore ask whether the first rotating state is the same object in motion or a different object altogether.

The wording is partly conventional, but the physical question is precise: how much energy is required to reach the first nonzero rotational state? If the gap is small, the excited object remains readily recognizable. If the gap is enormous, the excitation may change the object radically. An attempt to give an atom too much rotational excitation can throw off its electrons and destroy the atom rather than leave the same atom quietly rotating.

Now consider an electron, temporarily setting aside its known intrinsic spin. Can an electron be set rotating about an internal axis? A mathematical point has no visible internal structure to rotate, but experiment does not establish an electron's size by supplying a little rigid surface that can be watched turning. If the electron has an extremely small finite size, its first additional rotational excitation could simply require too much energy to reach. Because the energy rises as the size decreases, the required energy might be astronomical, perhaps even of order the Planck energy. This estimate is deliberately tentative.\footnote{A short passage immediately following this estimate is not intelligible enough to support a further substantive reconstruction.}

For an elementary particle, the magnitude of the spin is consequently treated as a fixed characteristic. Electrons do not come with a range of successively larger intrinsic spin magnitudes. If an otherwise electron-like excitation carried a different total spin, it would be classified as a different particle.

The fixed magnitude does not fix the direction. Imagine a particle called a \emph{spintron}, identified partly by the magnitude of its angular momentum. If its spin points along one axis, rotational invariance requires that the same kind of particle can also have its spin pointed along another axis. No spatial direction is intrinsically preferred. The amount of spin remains fixed while its orientation can change. Thus the problem has two apparently conflicting ingredients: a quantized magnitude and freely rotatable directions. The angular-momentum algebra will show how they coexist.

\begin{classroomqa}
\lecturetimestamp{00:12:25}
\audiencequestion{Why is energy required to spin a small particle?}
\lecturerresponse{For an object of mass \(M\) and characteristic radius \(R\), the moment of inertia is of order
\begin{equation}
I\sim cMR^2,
\end{equation}
where the numerical factor \(c\) depends on the shape and on how the mass is distributed. If material near the outside moves with speed \(v\), its ordinary momentum is of order \(Mv\), and angular momentum is momentum times distance:
\begin{equation}
L\sim MvR.
\end{equation}
Equivalently, \(L=I\omega\) with \(v\sim\omega R\). The rotational kinetic energy is
\begin{equation}
E_{\mathrm{rot}}=\frac{L^2}{2I}.
\end{equation}
At fixed \(M\) and fixed \(L\),
\begin{equation}
E_{\mathrm{rot}}\sim\frac{L^2}{2cMR^2}\propto\frac{1}{R^2}.
\end{equation}
The smaller the object, the more energy a given angular momentum costs. That inverse-square dependence is classical. The additional quantum-mechanical fact is that angular momentum changes in discrete steps, so the first allowed step cannot be made arbitrarily small.}
\end{classroomqa}

The mathematics about to be used is abstract and initially rather unintuitive. Its interest is that experimentally testable facts emerge from a small set of algebraic rules.

\section{Orbital Angular Momentum}

Begin with one particle moving in a circular orbit around an origin. Its angular momentum has magnitude equal to momentum times the perpendicular distance from the origin. In a plane one may call one sense of circulation positive and the reverse sense negative. The sign records the direction of the angular-momentum vector relative to the chosen axis.

This is not yet spin. The particle has angular momentum because it has ordinary momentum. Now place two particles at opposite ends of a diameter and let them move with the same sense of rotation. Their ordinary momenta cancel, so the center of mass is at rest, but their angular momenta point in the same direction and add. The system has zero total momentum and nonzero total angular momentum: a classical model of spin.

Let the position of a particle relative to the origin be
\begin{equation}
\vec r=(x,y,z)=(x_1,x_2,x_3),
\end{equation}
and let its momentum be
\begin{equation}
\vec p=(p_x,p_y,p_z)=(p_1,p_2,p_3).
\end{equation}
The angular momentum is the vector product
\begin{equation}
\vec L=\vec r\times\vec p.
\end{equation}
Reversing the factors reverses the vector: \(\vec p\times\vec r=-\vec r\times\vec p\). The standard convention is \(\vec r\times\vec p\), with its direction fixed by the right-hand rule.\footnote{\raggedright Lecture video frame \nolinkurl{lecture_07_figure_02.png}, 00:18:49.\par}

\begin{figure}[h]
\centering
\includegraphics[width=0.72\textwidth]{lecture_07_figure_02.png}
\par
\begin{tikzpicture}[scale=1.0]
\draw[->] (-0.6,0) -- (4.8,0) node[right] {$x$};
\draw[->] (0,-0.6) -- (0,3.8) node[above] {$y$};
\coordinate (P) at (3.2,2.3);
\draw[->,thick] (0,0) -- (P) node[midway,above left] {$\vec r$};
\draw[dashed] (P) -- (3.2,0);
\draw[dashed] (P) -- (0,2.3);
\draw[->,thick] (P) -- ++(-1.0,1.25) node[above] {$\vec p$};
\node at (4.15,2.75) {$\vec L=\vec r\times\vec p$};
\end{tikzpicture}
\caption{For a particle at position \(\vec r\) with momentum \(\vec p\), the angular momentum \(\vec L=\vec r\times\vec p\) is perpendicular to their plane.}
\lectureframe{00:18:49}
\end{figure}

In Cartesian components,
\begin{equation}
\begin{aligned}
L_x&=yp_z-zp_y,\qquad
L_y=zp_x-xp_z,\\
L_z&=xp_y-yp_x.
\end{aligned}
\end{equation}
It is enough to remember the first formula and cycle the labels
\[
x\longrightarrow y\longrightarrow z\longrightarrow x.
\]
Cycling both the component label and the labels on the right gives the next formula without changing the relative sign. The mnemonic is useful precisely because the sign is easy to reverse accidentally.

For a point particle, \(\vec r\times\vec p\) is orbital angular momentum about the chosen origin. For a composite system, add the angular momenta of all constituents as vectors:
\begin{equation}
\vec L_{\mathrm{tot}}=\sum_a \vec r_a\times\vec p_a.
\end{equation}
When this sum is evaluated in the center-of-mass frame, it can remain nonzero even though
\[
\sum_a\vec p_a=0.
\]
That remaining angular momentum is the spin of the composite system.

\section{Quantum Operators and Canonical Commutators}

In quantum mechanics, observable quantities are represented by operators. Position, momentum, and each component of angular momentum are therefore operators. Their detailed realization as differential operators is not needed here. The commutation relations contain enough information to derive the angular-momentum spectrum.

For two operators \(A\) and \(B\),
\begin{equation}
[A,B]=AB-BA.
\end{equation}
If two observables commute, they can be simultaneously specified. If they do not commute, quantum mechanics prevents their simultaneous sharp specification.

All three position coordinates can be measured together, so the position components commute:
\begin{equation}
[x_i,x_j]=0.
\end{equation}
In particular, every operator commutes with itself. It would make no sense for an observable to be incompatible with another measurement of that same observable.

The momentum components likewise commute:
\begin{equation}
[p_i,p_j]=0.
\end{equation}
A position component commutes with a momentum component in a different direction, but not with its matching momentum component. The three canonical statements can be written together as
\begin{equation}
[x_i,x_j]=0,\qquad
[p_i,p_j]=0,\qquad
[x_i,p_j]=i\hbar\,\delta_{ij}.
\end{equation}
Thus, for example,
\[
[x,p_y]=0,\qquad [x,p_x]=i\hbar,
\]
and reversing the order reverses the sign:
\begin{equation}
[p_j,x_i]=-i\hbar\,\delta_{ij}.
\end{equation}

Angular momentum and Planck's constant have the same dimensions. Angular momentum is distance times momentum, while the uncertainty relation has the dimensional form
\begin{equation}
\Delta x\,\Delta p\sim\hbar
\end{equation}
up to the familiar numerical factor in the precise inequality. It is therefore natural that angular momentum will be quantized in units of \(\hbar\), although the quantization has not yet been derived.

\section{Deriving the Angular-Momentum Algebra}

The first task is to compute the commutator of two angular-momentum components. From
\[
L_x=yp_z-zp_y,\qquad
L_y=zp_x-xp_z,
\]
one has
\begin{align}
[L_x,L_y]
&=[yp_z-zp_y,\;zp_x-xp_z] \\
&=[yp_z,zp_x]-[yp_z,xp_z] \notag\\
&\quad{}-[zp_y,zp_x]+[zp_y,xp_z].
\end{align}

\begin{classroomqa}
\lecturetimestamp{00:31:41}
\audiencequestion{Can the component formulas be remembered by cycling the symbols from left to right instead of cycling the three equations from top to bottom?}
\lecturerresponse{Yes. Either cycling mnemonic is acceptable. The only purpose of the mnemonic is to generate the components in the correct cyclic order while keeping track of the sign.}
\end{classroomqa}

A commutator can be understood by asking whether the factors in one product can be pushed through the factors in the other. In
\[
[yp_z,xp_z],
\]
the operator \(y\) commutes with both \(x\) and \(p_z\), while \(p_z\) commutes with both \(x\) and itself. There is no obstruction, so this term vanishes. The same reasoning gives
\[
[zp_y,zp_x]=0.
\]
Only two product pairs survive.

For the first one, the only obstruction is \(p_z\) passing through \(z\):
\begin{align}
[yp_z,zp_x]
&=yp_zzp_x-zp_xyp_z \\
&=y(p_zz-zp_z)p_x \\
&=y[p_z,z]p_x \\
&=-i\hbar\,yp_x.
\end{align}
The order of the canonical commutator matters:
\[
[p_z,z]=-[z,p_z]=-i\hbar.
\]

For the other surviving pair, the minus signs in the definitions of \(L_x\) and \(L_y\) multiply to a plus. The only obstruction is \(z\) passing through \(p_z\):
\begin{align}
[zp_y,xp_z]
&=zp_yxp_z-xp_zzp_y \\
&=x(zp_z-p_zz)p_y \\
&=x[z,p_z]p_y \\
&=i\hbar\,xp_y.
\end{align}
Combining the two contributions,
\begin{align}
[L_x,L_y]
&=-i\hbar\,yp_x+i\hbar\,xp_y \\
&=i\hbar(xp_y-yp_x) \\
&=i\hbar L_z.
\end{align}

This is the important simplification. Commuting two components of angular momentum produces the third component; no position or momentum operator remains. Cyclic permutation gives
\begin{equation}
\begin{aligned}
[L_x,L_y]&=i\hbar L_z,\qquad
[L_y,L_z]=i\hbar L_x,\\
[L_z,L_x]&=i\hbar L_y.
\end{aligned}
\end{equation}
Equivalently,
\begin{equation}
[L_i,L_j]=i\hbar\,\epsilon_{ijk}L_k.
\end{equation}

The algebra closes on the angular-momentum operators themselves. Once these relations have been obtained, the \(x_i\) and \(p_i\) used to derive them can be set aside. The same algebra applies when the total angular momentum is a sum over many constituents. In particular, it applies in the center-of-mass frame, where the remaining angular momentum is spin. The angular-momentum commutators can therefore be used abstractly without first constructing spin as the motion of little constituent parts.

The choice of \(x\), \(y\), and \(z\) axes is arbitrary. Rotate to any other mutually perpendicular set of inertial Cartesian axes, and the primed components obey the same relations:
\begin{equation}
[L'_x,L'_y]=i\hbar L'_z,
\end{equation}
together with its cyclic partners. The algebra is rotationally invariant, so physical predictions cannot depend on which inertial axes were chosen for the calculation.

\begin{classroomqa}
\lecturetimestamp{00:39:09}
\audiencequestion{Is spin angular momentum independent of translations as well as rotations of the coordinate axes?}
\lecturerresponse{Yes. Intrinsic spin is translation-independent. Orbital angular momentum about an origin can change when the origin is translated, but the spin remaining in the center-of-mass frame does not depend on that translation.}
\end{classroomqa}

Abstract algebra can generate endlessly many identities, most of them physically uninteresting. Its value appears when an identity becomes a measurable fact or reveals something about the operational meaning of measurement. Up to this point the construction is still mainly algebraic. The next step extracts the discrete structure of spin.

\begin{classroomqa}
\lecturetimestamp{00:40:42}
\audiencequestion{What if the chosen \(x\), \(y\), and \(z\) axes rotate along with the particles?}
\lecturerresponse{The discussion is restricted to inertial frames. A rotating coordinate system is noninertial and introduces effects such as centrifugal and Coriolis forces. Angular momentum is not invariant under a change to a co-rotating frame: a system that is visibly spinning in an inertial frame can appear stationary in a frame rotating with it. The rotational invariance used here concerns differently oriented inertial Cartesian axes, not axes that rotate in time.}
\end{classroomqa}

\section{Ladder Operators and the Quantization Axis}

Introduce the two combinations
\begin{equation}
L_+=L_x+iL_y,
\qquad
L_-=L_x-iL_y.
\end{equation}
They are analogous to the creation and annihilation operators of the harmonic oscillator. Their role is to raise and lower one chosen component of angular momentum.

There is nothing special about the \(z\)-axis, but one axis must be chosen if a component is to be specified. Take \(z\) as the quantization axis and write
\begin{equation}
L_z\lvert m\rangle=m\hbar\,\lvert m\rangle.
\end{equation}
Thus \(m\) is the measured \(z\)-component in units of \(\hbar\). It is often convenient during the algebra to set \(\hbar=1\); it will be retained explicitly here.

The commutator \([L_+,L_-]\) is proportional to \(L_z\), but it is not needed for the immediate argument. What matters first is the commutator of each ladder operator with \(L_z\). From
\[
[L_x,L_z]=-i\hbar L_y,
\qquad
[L_y,L_z]=i\hbar L_x,
\]
one obtains
\begin{align}
[L_+,L_z]
&=[L_x+iL_y,L_z] \\
&=[L_x,L_z]+i[L_y,L_z] \\
&=-i\hbar L_y+i(i\hbar L_x) \\
&=-\hbar(L_x+iL_y) \\
&=-\hbar L_+.
\end{align}
Changing the sign of the \(iL_y\) term gives
\begin{equation}
[L_-,L_z]=+\hbar L_-.
\end{equation}
Equivalently, with the order of each commutator reversed,
\begin{equation}
[L_z,L_+]=\hbar L_+,
\qquad
[L_z,L_-]=-\hbar L_-.
\end{equation}

Only one component of angular momentum can be sharp at a time. Any two distinct components fail to commute, so they cannot be simultaneously specified. The chosen sharp component defines the quantization axis; choosing \(z\) is a convenience, not a physical preference.

Now act with \(L_+\) on an \(L_z\) eigenstate. The commutator can be rewritten as
\begin{equation}
L_zL_+=L_+L_z+\hbar L_+.
\end{equation}
Therefore
\begin{align}
L_z\bigl(L_+\lvert m\rangle\bigr)
&=\bigl(L_+L_z+\hbar L_+\bigr)\lvert m\rangle \\
&=m\hbar L_+\lvert m\rangle+\hbar L_+\lvert m\rangle \\
&=(m+1)\hbar\,L_+\lvert m\rangle.
\end{align}
Provided it is nonzero, \(L_+\lvert m\rangle\) is an \(L_z\) eigenstate whose eigenvalue has increased by one unit of \(\hbar\). In the same way,
\begin{equation}
L_z\bigl(L_-\lvert m\rangle\bigr)
=(m-1)\hbar\,L_-\lvert m\rangle.
\end{equation}
Thus \(L_-\) lowers \(m\) by one.

This is already a nontrivial conclusion. Whatever the spectrum of \(L_z\) turns out to be, neighboring values are separated by integers:
\begin{equation}
\ldots,\ m-2,\ m-1,\ m,\ m+1,\ m+2,\ \ldots
\end{equation}
The endpoints have not yet been determined.

\section{Finite Multiplets and the Spin Label}

A fixed kind of particle has a fixed angular-momentum magnitude. Its angular momentum may be reoriented, but increasing its magnitude without limit would eventually produce a different object. There must therefore be a maximum possible \(z\)-component for that particle.

How can a ladder with a raising operator have a top? The harmonic oscillator supplies the model. Its annihilation operator gives zero when it reaches the bottom state. Here the raising operator must give zero when it reaches the top state:
\begin{equation}
L_+\lvert m_{\max}\rangle=0.
\end{equation}
The statement that \(L_+\) raises an eigenvalue applies only when the vector it produces is nonzero. The highest-weight state is precisely the exceptional state at which the raising sequence terminates.

Rotational invariance fixes the lower endpoint. If the angular momentum can point maximally along \(+z\), the same physical object can be rotated so that it points maximally along \(-z\). Hence
\begin{equation}
m_{\min}=-m_{\max}.
\end{equation}
The spectrum is symmetric about zero and has unit spacing. There are two possible families.

In the first family, zero occurs:
\begin{equation}
\begin{aligned}
m={}&-m_{\max},-m_{\max}+1,\ldots,-1,0,\\
&1,\ldots,m_{\max}-1,m_{\max}.
\end{aligned}
\end{equation}
Then \(m_{\max}\) is an integer, and the number of states is
\begin{equation}
2m_{\max}+1.
\end{equation}
The maximum value characterizes the particle and is called its spin. In conventional notation,
\begin{equation}
s\equiv m_{\max},
\qquad
\lvert s,m\rangle\equiv\lvert m\rangle
\end{equation}
within a fixed-spin multiplet.\footnote{\raggedright Lecture video frame \nolinkurl{lecture_07_figure_03.png}, 01:01:08.\par}

\begin{figure}[htbp]
\centering
\includegraphics[width=0.70\textwidth]{lecture_07_figure_03.png}
\caption{The highest-weight condition terminates the raising sequence, and the allowed \(m\) values form a finite unit-spaced ladder.}
\lectureframe{01:01:08}
\end{figure}

A spin-zero particle has only \(m=0\). A spin-one particle has
\[
m=1,0,-1,
\]
and a spin-two particle has five component values. The Higgs boson is an example of spin zero. The photon and the \(Z\) boson have spin one. A graviton, if present, has spin two. Composite objects can have spin three and much higher: nuclei can do so, and sufficiently macroscopic rotating objects such as basketballs and galaxies can carry arbitrarily large angular momentum. What is absent from the ordinary elementary-particle list is an elementary spin-three particle.

The other symmetric, unit-spaced family does not contain zero. It begins with
\[
m=\frac12,-\frac12
\]
and extends in unit steps:
\begin{equation}
\begin{aligned}
m={}&-m_{\max},-m_{\max}+1,\ldots,-\frac32,-\frac12,\\
&\frac12,\frac32,\ldots,m_{\max}.
\end{aligned}
\end{equation}
Here \(m_{\max}\) is a half-integer. The simplest member has spin \(\tfrac12\), followed by spin \(\tfrac32\), spin \(\tfrac52\), and so on. Angular-momentum multiplets therefore come in integer-spin and half-integer-spin families.

\begin{classroomqa}
\lecturetimestamp{01:03:46}
\audiencequestion{Does a half-integer-spin multiplet also have a maximum?}
\lecturerresponse{Yes. The ladder again has a highest value. For a spin-\(\tfrac32\) particle the maximum is \(m_{\max}=\tfrac32\); for a spin-\(\tfrac52\) particle it is \(m_{\max}=\tfrac52\). The difference is only that the maximum and every other value in the multiplet are half-integers.}
\end{classroomqa}

\begin{classroomqa}
\lecturetimestamp{01:04:04}
\audiencequestion{Are there no elementary particles of spin three or higher?}
\lecturerresponse{Within the relativistic quantum framework under discussion, the elementary-particle spins on the usual list are \(0,\tfrac12,1,\tfrac32,\) and \(2\). Composite objects are not restricted to this range and can have spin three or much larger spin.}
\end{classroomqa}

\begin{classroomqa}
\lecturetimestamp{01:04:50}
\audiencequestion{Do the negative entries on the ladder represent negative spin?}
\lecturerresponse{No. A negative \(m\) is a component value: it says that the angular momentum points oppositely relative to the chosen axis. The spin of the particle is the nonnegative maximum \(s=m_{\max}\).}
\end{classroomqa}

\begin{classroomqa}
\lecturetimestamp{01:05:10}
\audiencequestion{If the component values are discrete, what happens to all the orientations between them?}
\lecturerresponse{There are not merely two or three possible spatial orientations. Intermediate orientations are represented by quantum superpositions of the \(m\)-basis states. The spin-\(\tfrac12\) construction below makes this explicit.}
\end{classroomqa}

\begin{classroomqa}
\lecturetimestamp{01:05:42}
\audiencequestion{What about corkscrew motion or circularly polarized waves?}
\lecturerresponse{A circularly polarized electromagnetic wave is built from circularly polarized photons, and those photons carry angular momentum along the propagation axis. The word ``corkscrew'' combines rotation with forward motion, so it is not the right description for an object being considered at rest. For a moving particle, the relevant quantity is helicity: the orientation of its spin relative to its direction of motion.}
\end{classroomqa}

\begin{classroomqa}
\lecturetimestamp{01:07:10}
\audiencequestion{Is any known elementary particle spin \(\tfrac32\)?}
\lecturerresponse{No spin-\(\tfrac32\) elementary particle has been directly confirmed. Supersymmetry conjectures a spin-\(\tfrac32\) partner of the graviton called the gravitino, but it has not been directly detected.}
\end{classroomqa}

\begin{classroomqa}
\lecturetimestamp{01:08:01}
\audiencequestion{Have the angular-momentum operators stopped acting on the original wavefunction or ket space?}
\lecturerresponse{No original spin space had yet been specified. The algebra is being used to determine that space. For fixed spin \(s\), it has basis vectors labeled by the allowed values
\[
\lvert s,m\rangle,\qquad m=-s,-s+1,\ldots,s.
\]
This intrinsic spin space is distinct from the ordinary space of spatial wavefunctions on which orbital angular momentum acts as \(\mathbf r\times\mathbf p\), although both realize the same angular-momentum algebra.}
\end{classroomqa}

\section{The Magnitude of Angular Momentum}

\begin{classroomqa}
\lecturetimestamp{01:09:32}
\audiencequestion{Is the magnitude of the angular-momentum vector observable?}
\lecturerresponse{Yes. It is most convenient to use its square,
\[
L^2=L_x^2+L_y^2+L_z^2.
\]
Its value on the highest-weight state is not simply \(s^2\hbar^2\). Even when \(L_z\) is sharp and maximal, the noncommuting transverse components cannot both be sharply zero, so they supply an additional quantum contribution.}
\end{classroomqa}

To calculate that contribution, multiply the ladder operators in the order for which the raising operator acts first on a ket:
\begin{align}
L_-L_+
&=(L_x-iL_y)(L_x+iL_y) \\
&=L_x^2+L_y^2+iL_xL_y-iL_yL_x \\
&=L_x^2+L_y^2+i[L_x,L_y] \\
&=L_x^2+L_y^2-\hbar L_z.
\end{align}
The final step uses \([L_x,L_y]=i\hbar L_z\). Therefore
\begin{equation}
L_x^2+L_y^2=L_-L_++\hbar L_z,
\end{equation}
and hence
\begin{equation}
L^2=L_-L_++L_z^2+\hbar L_z.
\end{equation}

Apply this identity to the highest-weight state \(\lvert s,s\rangle=\lvert m_{\max}\rangle\):
\begin{align}
L^2\lvert s,s\rangle
&=\left(L_-L_++L_z^2+\hbar L_z\right)\lvert s,s\rangle \\
&=L_-L_+\lvert s,s\rangle
+s^2\hbar^2\lvert s,s\rangle
+s\hbar^2\lvert s,s\rangle.
\end{align}
The first term vanishes because \(L_+\lvert s,s\rangle=0\). Thus
\begin{equation}
L^2\lvert s,s\rangle
=s(s+1)\hbar^2\lvert s,s\rangle.
\end{equation}
The extra \(+s\) is the quantum correction associated with the unavoidable transverse fluctuations. Classically one would have expected only \(s^2\hbar^2\).

The first few values make the distinction concrete:
\begin{align}
s=0 &: \qquad L^2=0,\\
s=1 &: \qquad L^2=2\hbar^2,\\
s=\frac12 &: \qquad L^2=\frac34\hbar^2.
\end{align}
For \(s=1000\), the dimensionless factor is \(1000\times1001\), which differs little from \(1000^2\) on the scale of such a large angular momentum. In general,
\begin{equation}
s(s+1)=s^2\left(1+\frac1s\right),
\end{equation}
so the relative quantum correction becomes negligible for very large \(s\). That is how the classical \(s^2\) result emerges.

Every \(m\)-state in a fixed spin multiplet has this same eigenvalue of \(L^2\). Changing \(m\) changes the component along the chosen axis, not the total magnitude. This is the quantum version of rotating an angular-momentum vector without changing its length.

\begin{classroomqa}
\lecturetimestamp{01:16:43}
\audiencequestion{Can \(L^2\) and the \(z\)-component be known simultaneously?}
\lecturerresponse{Yes. The total square commutes with every angular-momentum component:
\[
[L^2,L_i]=0,\qquad i=x,y,z.
\]
Thus \(L^2\) and one chosen component, such as \(L_z\), can be specified together even though two distinct components cannot. Every electron, for example, has \(L^2=\tfrac34\hbar^2\), independent of its orientation.}
\end{classroomqa}

\begin{classroomqa}
\lecturetimestamp{01:17:41}
\audiencequestion{Is \(L^2\) zero for spin one, or only for spin zero?}
\lecturerresponse{Only spin zero has \(L^2=0\). Spin one gives
\[
L^2=1(1+1)\hbar^2=2\hbar^2,
\]
while spin three gives
\[
L^2=3(3+1)\hbar^2=12\hbar^2.
\]}
\end{classroomqa}

\begin{classroomqa}
\lecturetimestamp{01:18:13}
\audiencequestion{Why does every \(m\)-level in the multiplet have the same \(L^2\) eigenvalue?}
\lecturerresponse{Start from the highest-weight state and lower it. Since \(L^2\) commutes with every component, it also commutes with \(L_-\):
\[
[L^2,L_-]=0.
\]
If
\[
L^2\lvert s,s\rangle=s(s+1)\hbar^2\lvert s,s\rangle,
\]
then
\begin{align}
L^2\bigl(L_-\lvert s,s\rangle\bigr)
&=L_-L^2\lvert s,s\rangle \\
&=s(s+1)\hbar^2\,L_-\lvert s,s\rangle.
\end{align}
Thus the next state down has the same \(L^2\) eigenvalue. Repeating the argument carries that eigenvalue all the way down the ladder.}
\end{classroomqa}

The discrete \(m\)-values therefore do not mean that only a discrete set of spatial orientations exists. Rotational invariance still requires states corresponding to rotated directions. For spin \(\tfrac12\), those directions will appear as superpositions of the two \(L_z\) basis states.

\section{Spin-\texorpdfstring{\(\tfrac12\)}{1/2} States and Orientation}

Rotational invariance requires more than the discrete \(L_z\) eigenstates. A spin may point along any axis, even though measurement of its component along a chosen axis gives only the allowed discrete values. The directions between the \(z\)-axis eigenstates are represented by superpositions.

For a spin-\(\tfrac12\) particle, restrict attention to the spin degrees of freedom. An electron also has orbital states, so it certainly does not have only two states altogether; its spin space alone is two-dimensional. Choose the \(z\)-axis basis
\begin{equation}
L_z\lvert+\rangle=\frac{\hbar}{2}\lvert+\rangle,
\qquad
L_z\lvert-\rangle=-\frac{\hbar}{2}\lvert-\rangle.
\end{equation}
A general spin state is
\begin{equation}
\lvert\psi\rangle
=
\alpha\lvert+\rangle+\beta\lvert-\rangle,
\end{equation}
where \(\alpha\) and \(\beta\) are complex numbers. Measurement of \(L_z\) gives \(+\hbar/2\) with probability \(|\alpha|^2\) and \(-\hbar/2\) with probability \(|\beta|^2\), so normalization requires
\begin{equation}
|\alpha|^2+|\beta|^2=1.
\end{equation}

A common phase has no physical effect:
\begin{equation}
(\alpha,\beta)\longrightarrow
e^{i\chi}(\alpha,\beta).
\end{equation}
The two complex amplitudes begin with four real parameters. Normalization removes one, and the irrelevant overall phase removes another, leaving two physical real parameters. That is exactly the number needed to specify a direction by its polar and azimuthal angles. Varying \(\alpha\) and \(\beta\) rotates the spin orientation.

A measured component and its expectation value are different things. The expectation value is an ensemble average. Prepare many electrons in exactly the same state, measure \(L_z\) once on each electron, add the results, and divide by the number of experiments. Every individual result is still either \(+\hbar/2\) or \(-\hbar/2\), but the average may take any value between them:
\begin{equation}
\langle L_z\rangle
=
\frac{\hbar}{2}\bigl(|\alpha|^2-|\beta|^2\bigr).
\end{equation}

Equal up and down amplitudes give a spin directed along one of the \(x\)-axis orientations. Normalized representatives are\footnote{The coefficient ratios fix the directions. The factors \(1/\sqrt2\) supply the normalization required by \( |\alpha|^2+|\beta|^2=1\). Relative phases \(+i\) and \(-i\) give opposite \(y\)-axis orientations; assigning either sign to positive \(y\) depends on the axis convention.}
\begin{equation}
\lvert x;+\rangle
=
\frac{\lvert+\rangle+\lvert-\rangle}{\sqrt2},
\qquad
\lvert x;-\rangle
=
\frac{\lvert+\rangle-\lvert-\rangle}{\sqrt2}.
\end{equation}
For either state, measurement along \(z\) gives the two possible outcomes with equal probability, and therefore
\begin{equation}
\langle L_z\rangle
=
\frac12\left(\frac{\hbar}{2}\right)
+
\frac12\left(-\frac{\hbar}{2}\right)
=0.
\end{equation}
The plus and minus relative signs give opposite directions along the \(x\)-axis. Similarly,
\begin{equation}
\frac{\lvert+\rangle+i\lvert-\rangle}{\sqrt2},
\qquad
\frac{\lvert+\rangle-i\lvert-\rangle}{\sqrt2}
\end{equation}
give the two opposite \(y\)-axis orientations. General normalized values of \(\alpha\) and \(\beta\) describe arbitrary axes. To say that a spin points along a particular axis means that measurement of the component along that axis gives a definite result. Relative to a different axis, the same state is generally a superposition.

The individual measurement outcomes remain quantized. It is their ensemble averages that vary continuously and reproduce the familiar geometry of a classical vector.

\section{Spin, Statistics, and Composite Particles}

Mass, electric charge, and spin are among the principal properties used to characterize a particle. Spin is also correlated with statistics. The spin-statistics relation is a theorem of relativistic quantum mechanics, although it will not be proved here:
\begin{equation}
\begin{aligned}
\text{half-integer spin} &\quad\Longleftrightarrow\quad \text{fermion},\\
\text{integer spin} &\quad\Longleftrightarrow\quad \text{boson}.
\end{aligned}
\end{equation}
Two identical fermions cannot occupy the same complete quantum state. Identical bosons may occupy the same state, and large numbers of them in one state can behave collectively like a classical wave. Thus spin \(\tfrac12,\tfrac32,\tfrac52,\ldots\) particles are fermions, while spin \(0,1,2,\ldots\) particles are bosons.

The rule extends immediately to composites. Adding two half-integer angular momenta gives an integer angular momentum. More generally, an even number of fermionic constituents makes a bosonic composite, whereas an odd number makes a fermionic composite. Adding bosonic constituents does not change this parity rule.

Ordinary hydrogen contains a spin-\(\tfrac12\) proton and a spin-\(\tfrac12\) electron. It therefore contains an even number of fermions and is bosonic as a whole. A deuteron contains a proton and a neutron, both fermions, and is also a boson. A deuterium atom adds an electron to the deuteron, giving three fermionic constituents at this level of description, so deuterium is a fermion.

Particles and antiparticles have the same mass and spin and opposite electric charge. Positronium is an electron and a positron bound to one another. It contains two fermions and is therefore bosonic.

\begin{classroomqa}
\lecturetimestamp{01:33:08}
\audiencequestion{Do bosons pass through things, and is superconductivity related to paired electrons behaving as bosons?}
\lecturerresponse{The relevant property is not simply an ability to pass through matter; it is the ability of bosons to occupy the same state. Superfluid helium gives the simpler example. An ordinary helium-\(4\) atom contains two protons, two neutrons, and two electrons, an even number of fermionic constituents, so the atom is a boson. Many helium atoms can occupy one quantum state and move collectively. A superfluid can then be pictured, approximately, as a classical wave formed by piling many atoms into the same state. That picture is not exact, but it captures the point needed here. An electron in a metal is a fermion and cannot share a complete state with another electron. Pairs of electrons, however, are bosonic composites, which is the connection to superconductivity. This raises a paradox: if the constituent electrons cannot occupy the same state, how can several composite pairs occupy one state? The resolution requires a more careful account of composite states and is deferred.}
\end{classroomqa}

\begin{classroomqa}
\lecturetimestamp{01:35:46}
\audiencequestion{Since protons and neutrons are composite, why do they have half-integer spin?}
\lecturerresponse{Each is made from an odd number of spin-\(\tfrac12\) quarks. Three fermionic constituents make a fermionic composite, so protons and neutrons can have half-integer spin.}
\end{classroomqa}

\section{Pauli Exclusion and Atomic Orbitals}

The periodic table provides the immediate motivation for the Pauli exclusion principle. The rule is not that two electrons can never have the same orbital motion. It is that two electrons cannot occupy the same complete quantum state.

This distinction matters even in the simplest atoms. Starting from hydrogen, increase the nuclear charge and add a second electron. The two electrons may occupy the same spatial orbital, but their spin states must differ. One may be spin up and the other spin down with respect to the chosen axis. When two electrons occupy one orbital, their spins combine into the spin-singlet state, so their total spin cancels. No more than two electrons may occupy a single spatial orbital because there are only two independent spin states available.

\begin{classroomqa}
\lecturetimestamp{01:37:53}
\audiencequestion{Can the opposite spins of two electrons in one orbital be pictured as meshed gears?}
\lecturerresponse{As a rough picture, yes: meshed gears turn in opposite senses. It is only an analogy for the opposite spin directions, not a literal mechanical model of the electrons.}
\end{classroomqa}

The word ``orbital'' can be used at two different levels. A state of orbital motion, in the broad sense, is the state obtained by solving the Schrödinger equation while ignoring spin. It includes more information than merely which shell the electron occupies.

\begin{classroomqa}
\lecturetimestamp{01:38:09}
\audiencequestion{Does an orbital mean more than merely which shell the electron occupies?}
\lecturerresponse{Yes. Here the orbital motion means the electron's spin-ignored Schrödinger state. It includes the principal quantum number \(n\), which carries radial information and roughly sets the electron's distance from the nucleus, as well as the orbital-angular-momentum quantum numbers. In the spectroscopic notation used by chemists, \(S\), \(P\), \(D\), and \(F\) correspond respectively to orbital angular momentum \(l=0,1,2,3\).}
\end{classroomqa}

Thus the principal quantum number and the orbital angular momentum describe different aspects of the motion. The \(S,P,D,F\) labels refer to the value of \(l\), while \(n\) supplies separate radial information. Pauli's rule is that no two electrons may have all of their quantum numbers the same, including the spin quantum number.

\begin{classroomqa}
\lecturetimestamp{01:40:42}
\audiencequestion{Can the quantum numbers determine how many electrons may occupy each shell in an atom?}
\lecturerresponse{Yes. They determine the maximum possible number, though not how many electrons a particular atom actually contains in that shell. The general shell-counting problem will not be pursued here.}
\end{classroomqa}

For a fixed orbital angular momentum \(l\), the component quantum number has
\begin{equation}
2l+1
\end{equation}
possible values. This is the same angular-momentum counting already used for spin. An \(S\) state has \(l=0\), so it has one orbital state. A \(P\) state has \(l=1\), so it has three orbital states. A \(D\) state has \(l=2\), so it has five.

Spin doubles each orbital count because each spatial orbital can carry one electron in each of the two opposite spin states.

\begin{classroomqa}
\lecturetimestamp{01:41:54}
\audiencequestion{Does the \(S\) state hold two electrons, and do the three \(P\) states hold six more?}
\lecturerresponse{Yes. The single \(S\) orbital state holds at most two electrons. The three \(P\) orbital states hold two each, giving six more. The \(S\) and \(P\) states under discussion therefore accommodate \(2+6=8\) electrons in total.}
\end{classroomqa}

The same angular-momentum structure that organizes atomic orbitals appears again, with some differences, when quarks move in bound systems and orbit one another.
